\section{The integer dynamics and admissible forcing}\label{sec:foundations}

We use the shortcut Collatz map on the signed integers,
\begin{equation}\label{eq:shortcut}
 T(n)=\begin{cases}n/2,&n\equiv0\pmod2,\\
 (3n+1)/2,&n\equiv1\pmod2.
 \end{cases}
\end{equation}
The divisions in this formula are exact. A forcing word is a nonempty
finite list $d=(d_0,\ldots,d_{p-1})$ of integers, repeated periodically.
The augmented state and its update are
\begin{equation}\label{eq:physical-forcing}
 (n,i)\in\Z\times(\Z/p\Z),\qquad
 F_d(n,i)=\bigl(T(n)+d_i,i+1\bmod p\bigr).
\end{equation}
Write $F_{d,i}^{(k)}(n)$ for the scalar value after $k$ such updates
starting at phase $i$. The impulse is applied before the phase advances.
Zero continues to evolve by this same rule; it is not an absorbing state.
We allow a displayed period $p$ to be a multiple of the least, or
\emph{primitive}, forcing period. When a bound is expressed relative to
the primitive period, that qualification is explicit.

\subsection{A trajectory includes its clock}

The smallest nontrivial example used repeatedly below is
\begin{equation}\label{eq:three-example-word}
 d=(1,0,-1).
\end{equation}
Phase zero applies $T$ and then adds one; phase one applies $T$;
phase two applies $T$ and then subtracts one. For example, a launch
of seed $5$ at phase zero begins
\[
 (5,0)\longmapsto(9,1)\longmapsto(14,2)\longmapsto(6,0).
\]
The first update is $T(5)+1=8+1=9$. In particular, the impulse
is not inserted into the argument of $T$: $T(5+1)=3$ would give
a different recurrence.

The following comparison keeps the seed fixed and changes only the
clock position. Entries are scalar values; the forced phase at time
$k$ is the launch phase plus $k$ modulo three.
\begin{center}
\begin{tabular}{rrrrr}\toprule
$k$ & Ordinary $T^k(5)$ & Launch at $0$ & Launch at $1$ & Launch at $2$\\\midrule
0&5&5&5&5\\
1&8&9&8&7\\
2&4&14&3&12\\
3&2&6&6&6\\
4&1&4&3&2\\
5&2&2&4&2\\
6&1&0&3&1\\\bottomrule
\end{tabular}
\end{center}
At time three, all three forced scalar values happen to equal $6$.
Their phases are different, so their next values are respectively
$4$, $3$, and $2$. Thus scalar equality does not by itself give a
merger. Equality of the augmented states does: the update is a
function of precisely that pair. This distinction will be used in
both the inverse constructions and the later rank reductions.

The phase-zero launch also illustrates continuation through zero.
Its next three updates after reaching $(0,0)$ are
\[
 (0,0)\longmapsto(1,1)\longmapsto(2,2)\longmapsto(0,0).
\]
By contrast, a launch from $(0,1)$ follows
$(0,1)\mapsto(0,2)\mapsto(-1,0)$. A statement about every positive
seed can therefore coexist with restrictions on which zero states
are allowed. The complete cycle geometry for this word is displayed
in Figure~\ref{fig:phase-cycles}; the present calculation explains
why the phase labels in that figure are part of the mathematics.

\begin{definition}\label{def:admissible}
A forcing word is \emph{conservative} if $\sum_{i=0}^{p-1}d_i=0$.
It is \emph{admissible} if, for every $n\ge1$, every launch phase $i$,
and every $k\ge0$, one has $F_{d,i}^{(k)}(n)\ge0$.
A nonnegative state is \emph{safe} if all of its subsequent scalar
values are nonnegative.
\end{definition}

Conservation is a condition on the additive impulses. It is not a
conservation law for the integer state. Admissibility quantifies over
every positive seed at every phase; safety of one observed trajectory
does not imply it.

\begin{theorem}[Admissible integer words]\label{thm:admissibility}
A periodic integer word is admissible if and only if every impulse is
at least $-1$ and each pair of successive negative impulses in cyclic
order has a positive impulse between them. For a word with a single
negative impulse, the intervening arc means one full circuit back to
that impulse. The separation requirement is vacuous if there are no
negative impulses.

Consequently, a conservative word is admissible if and only if it is
identically zero or its nonzero entries alternate cyclically between
$+1$ and $-1$.
\end{theorem}
\begin{proof}
Launching seed $2$ at an impulse $d_i\le-2$ gives
$T(2)+d_i=1+d_i<0$, so every negative impulse of an admissible word
equals $-1$. If two successive negative impulses have only zeros
between them, seed $2$ at the first follows $2\mapsto0$, stays zero
until the second, and then becomes $-1$. This proves necessity,
including a single negative impulse repeated on the next lap.

Conversely, $T(n)\ge1$ when $n\ge1$. A negative impulse can therefore
create zero but cannot create a negative state from a positive input.
After such a zero, zero impulses leave it at zero and a positive
impulse restores a positive value. The separation condition guarantees
this restoration before the next negative impulse. Thus every
positive launch at every phase stays nonnegative.

Suppose now that the word is nonzero and conservative. It has both
positive and negative entries. Let $a$ and $b$ count these entries,
respectively. Separation gives $a\ge b$. Since every negative entry
is $-1$, conservation gives
$\sum_{d_i>0}d_i=b$, whereas integrality gives
$\sum_{d_i>0}d_i\ge a$. Hence $a=b$, all positive entries are $1$,
and exactly one positive entry separates successive negative entries.
The converse follows from the first assertion.
\end{proof}

The second generator used below is the translated Collatz map
$U=P\circ T\circ P^{-1}$, where $P(x)=x+1$. It is exactly the map
$E$ of Monks et al.~\cite[\S~3]{monks-et-al-2013}.
The forcing-specific assertions here identify the binary height, the
safe domain, and the phase-dependent coordinate that selects between
these two generators.

\begin{proposition}[Binary height and positive coordinates]
\label{prop:height-conjugacy}
Let $d$ be a nonzero conservative admissible word. There is a unique
periodic height word $h_i\in\{0,1\}$ such that
\begin{equation}\label{eq:height-difference}
 d_i=h_{i+1}-h_i.
\end{equation}
The height and impulse words have the same primitive period. Their
safe nonnegative states are exactly
\begin{equation}\label{eq:safe-domain}
 \{(n,i):n\ge h_i\}.
\end{equation}
On this domain, $z=n+1-h_i$ is a positive integer and conjugates the
physical dynamics to
\begin{equation}\label{eq:positive-switch}
 z'=\begin{cases}T(z),&h_i=1,\\ U(z),&h_i=0,
 \end{cases}
 \qquad
 U(z)=T(z-1)+1=
 \begin{cases}3z/2,&z\text{ even},\\(z+1)/2,&z\text{ odd}.
 \end{cases}
\end{equation}
\end{proposition}
\begin{proof}
Set the height to zero just before a positive impulse, increase it to
one there, leave it constant at zero impulses, and decrease it to zero
at the next negative impulse. Alternation makes this definition
consistent around the clock. The value immediately before a positive
impulse must be zero in any binary realization, which proves uniqueness.

If the differences have period $q$, the quantity $h_{i+q}-h_i$ is
independent of $i$. Repeating this translation around the finite
clock shows that this constant is zero. Thus every period of the
differences is a period of the height, and the reverse implication is
immediate from \eqref{eq:height-difference}.

When $n\ge h_i$, one has $T(n)\ge h_i$. Consequently
$n'=T(n)+h_{i+1}-h_i\ge h_{i+1}$, proving invariance of
\eqref{eq:safe-domain}. The only nonnegative states outside it have
$n=0$ and $h_i=1$. From such a state the height remains one through
zero impulses until the next negative impulse, which sends zero to
$-1$. Such a negative impulse exists because the height is
nonconstant. This proves the safe-state characterization. Finally,
\[
 z'=n'+1-h_{i+1}=T(n)+1-h_i.
\]
Substitution of $n=z-1+h_i$ gives both cases of
\eqref{eq:positive-switch}.
\end{proof}

For zero forcing either constant height is possible, so the uniqueness
assertion deliberately excludes that word. We use $h_i=0$ if a
height is needed for zero forcing. In a nonzero conservative admissible word, zero
is a safe initial value precisely at height-zero phases. Every
positive initial value is safe at every phase.

The alternative coordinate $u=n-h_i=z-1$ is nonnegative. Its two
updates are $u'=T(u+h_i)-h_i$. We use $z$ in rank comparisons and
$n$ in statements about physical seeds. A divisibility restriction on
$z$ need not be the same restriction on $n$.

\subsection{Calculating with the height coordinate}

For \eqref{eq:three-example-word}, the unique binary height is
$h=(0,1,1)$. Taking cyclic differences gives
$(1-0,1-1,0-1)=(1,0,-1)$, and the dictionary between the two
coordinate systems is
\begin{center}
\begin{tabular}{rrrrll}\toprule
Phase $i$ & $d_i$ & $h_i$ & Safe physical values & Positive coordinate & Update\\\midrule
0&$+1$&0&$n\ge0$&$z=n+1$&$U$\\
1&0&1&$n\ge1$&$z=n$&$T$\\
2&$-1$&1&$n\ge1$&$z=n$&$T$\\\bottomrule
\end{tabular}
\end{center}
The chronological positive-coordinate word is consequently $UTT$.
The phase-zero path just calculated becomes
\begin{equation}\label{eq:height-worked-path}
 6\xrightarrow{U}9\xrightarrow{T}14\xrightarrow{T}7.
\end{equation}
The physical endpoint is $6$, while the positive-coordinate endpoint
is $7$, because the return is to height zero. This is a useful check
on a return-map calculation: comparison of a physical input with a
normalized output would mix two different quantities.

The formula for $U$ also explains the change in parity roles. If
$z=2m$, then $z-1$ is odd and
\[
 U(2m)=T(2m-1)+1=(3(2m-1)+1)/2+1=3m.
\]
If $z=2m+1$, then $U(2m+1)=T(2m)+1=m+1$.
For $T$, even inputs shrink and odd inputs grow; for $U$, those
roles are reversed. A fixed height word prescribes which of these
two maps is used at each time. It does not let an individual orbit
select the shrinking branch whenever convenient.

The safe-domain condition can now be read without subtraction at
the boundary: $n\ge h_i$ is exactly $z\ge1$. In particular, the
safe physical zero $(0,0)$ corresponds to $z=1$, whereas the unsafe
zero states at heights one would correspond to $z=0$. Positivity
of $z$ is preserved by both displayed formulas for $T$ and $U$.
This supplies a second way to follow the invariant domain, while
the identity in Proposition~\ref{prop:height-conjugacy} proves that
it describes the original physical dynamics.

There are therefore three compatible views of one admissible orbit:
physical value $n$ with its clock, the height chart $h_i$, and positive
value $z$. The first makes the forcing explicit, the second records
where safe zero states occur, and the third makes every update a
positive $T$ or $U$ step. We state the coordinate whenever a rank,
congruence, or base set depends on it.

\subsection{Affine branches and phase conventions}

For ordinary Collatz, the affine iterate formula and the correspondence
between finite parity words and residue classes are classical
\cite[Theorems~1.1--1.2]{terras-1976}; their $2$-adic coding is developed
in \cite[\S~1 and Appendix~B]{bernstein-lagarias-1996}.
The following calculation retains the residue-lifting argument and
includes the prescribed forcing in the affine intercept.

\begin{proposition}[Finite parity words]\label{prop:affine-parity}
Fix a forcing word, a launch phase $i$, and a finite binary word
$s_0,\ldots,s_{L-1}$. There is exactly one residue class modulo $2^L$
of integer seeds whose first $L$ physical input parities are this word.
On that class,
\begin{equation}\label{eq:affine-word}
 F_{d,i}^{(L)}(n)=\frac{3^{s_0+\cdots+s_{L-1}}n+B_L}{2^L},
 \qquad B_0=0,\quad
 B_{j+1}=3^{s_j}B_j+2^j(s_j+2d_{i+j}).
\end{equation}
Indices of $d$ are read modulo $p$.
\end{proposition}
\begin{proof}
For a step with input parity $s_j$,
$2n_{j+1}=3^{s_j}n_j+s_j+2d_{i+j}$.
Induction gives the displayed recurrence. For the residue assertion,
suppose a parity prefix of length $j$ specifies one residue modulo
$2^j$. Its two lifts modulo $2^{j+1}$ have seeds differing by $2^j$.
Their first $j$ parities agree, and their values after $j$ steps
differ by $3^{s_0+\cdots+s_{j-1}}$, which is odd. Exactly one lift
has either prescribed next parity. Induction starting at $j=0$
proves existence and uniqueness. Adding a multiple of $2^L$
preserves the entire prefix, so the identity holds on the whole class.
\end{proof}

This proposition is a finite itinerary statement. It does not assert
that an arbitrary infinite parity word has a positive integer seed.
Bernstein's explicit inverse coding \cite{bernstein-1994}, expressed
for the shortcut clock in \cite[\S~1]{bernstein-lagarias-1996}, takes
values in $\Z_2$; this ambient space does not enforce positive integer
realization.
The affine formula also explains why $\sum d_i=0$ need not cancel the
additive term $B_L$: later parity branches weight different impulses
differently.

\paragraph{An entire residue class from one parity calculation.}
Take phase zero in \eqref{eq:three-example-word} and prescribe the
three physical input parities $(1,1,0)$. The first odd input has
the form $n=2a+1$, and its first image is $3a+3$. Requiring that
image to be odd forces $a$ to be even, so $n=4b+1$. The next image
is $9b+5$; requiring it to be even forces $b$ to be odd. Thus the
successive parity constraints select
\[
 n\equiv1\pmod2,\qquad n\equiv1\pmod4,\qquad
 n\equiv5\pmod8.
\]
Writing $n=8q+5$ now gives a literal calculation on the whole class:
\begin{equation}\label{eq:worked-affine-cylinder}
 (8q+5,0)\longmapsto(12q+9,1)
 \longmapsto(18q+14,2)\longmapsto(9q+6,0).
\end{equation}
Every asserted parity is independent of $q$. For $q\ge0$ all
intermediate values are positive. The affine recurrence gives
$B_0=0$, $B_1=3$, $B_2=11$, $B_3=3$, so this branch of the
three-step return is
\[
 M(n)=\frac{9n+3}{8},\qquad n\equiv5\pmod8.
\]
For example, $5\mapsto9\mapsto14\mapsto6$ and
$13\mapsto21\mapsto32\mapsto15$. The formula describes infinitely
many seeds because the parity conditions and arithmetic were checked
with the free integer $q$, rather than inferred from those two seeds.
The conservative impulse sum is zero, but the branch has multiplier
$9/8$ and bias $3/8$.

To use this branch a second time, its endpoint must satisfy the same
guard. The condition $9q+6\equiv5\pmod8$ is $q\equiv7\pmod8$,
so two consecutive copies of the itinerary require
$n\equiv61\pmod{64}$. For instance,
\[
 61\longmapsto93\longmapsto140\longmapsto69
 \longmapsto105\longmapsto158\longmapsto78.
\]
More generally, induction on $k$ gives
\begin{equation}\label{eq:repeated-parity-cylinder}
 M^k(n)=\frac{9^k(n+3)}{8^k}-3,
 \qquad n\equiv-3\pmod{8^k},
\end{equation}
for the itinerary $(1,1,0)$ repeated $k$ times. The congruence is
both necessary and sufficient: the first guarded update sends
$8^kq-3$ to $9\cdot8^{k-1}q-3$, and the same argument applies
to the remaining copies. Every finite such itinerary has positive
realizers, for example $n=8^k-3$. An integer satisfying the guard
for every $k$ would have $8^k\mid n+3$ for all $k$, hence would
equal $-3$. The finite positive realizers therefore cannot be
replaced by one positive realizer of the infinite repeated word.
This explicit example separates a finite parity cylinder from a
claim about an infinite integer orbit.

Rotation has an exact augmented-state interpretation. If
$\widetilde d_j=d_{j+r}$, the map $(n,j)\mapsto(n,j+r)$ conjugates
$F_{\widetilde d}$ to $F_d$. This does not say that the same scalar
seed reaches the same cycle when launched at different phases of
one word.

An augmented cycle returns both the scalar value and the phase. Its
length is a multiple of the displayed clock length. If the clock has
primitive period $p$, any periodic scalar response has period a
multiple of $p$ as well: a scalar period $q$ makes
$d_k=n_{k+1}-T(n_k)$ periodic with period $q$, so $p$ divides $q$.
Thus the least scalar and augmented response periods agree on the
primitive clock. Repeating a shorter word to obtain a larger displayed
clock changes the augmented labeling but not the scalar recurrence.

\begin{proposition}[Nonnegative orbit dichotomy]
\label{prop:nonnegative-dichotomy}
For a fixed deterministic periodic forcing, a nonnegative orbit is
either eventually periodic or tends to $+\infty$. In particular,
boundedness, eventual periodicity, and failure of unboundedness are
equivalent for such an orbit.
\end{proposition}
\begin{proof}
A repeated augmented state forces a periodic future. If no state
repeats, then for each $M\ge0$ the orbit visits the finite set
$\{0,\ldots,M\}\times(\Z/p\Z)$ only finitely often. After its last
such visit every scalar value exceeds $M$. This is precisely escape
to $+\infty$. Conversely, every eventually periodic orbit is bounded.
\end{proof}

The proposition uses nonnegativity and a finite forcing clock. It
does not turn a finite trajectory with large peaks into a proof of
divergence, and it is not a statement about arbitrary signed or
aperiodically controlled trajectories.
